% TC-0039 fragment; intended for the C99 test report.
\def\TCID{TC-0039}
\def\TCTitle{C99 narrow numeric text interchange}
\def\TCPurpose{Verify hexadecimal floating syntax, special values, signed
zero, rounding, range behavior, and narrow formatted-I/O integration.}
\def\TCPriority{\TCPriorityHigh}
\def\TCRequirementRef{REQ-0039}
\def\TCDesignRef{ISO/IEC 9899:1999 7.19.6 and 7.20.1.3; documented
binary32/binary64 model.}
\def\TCFeatureRef{\texttt{src/stdlib.c}; internal numeric-text code;
\texttt{src/stdio\_format.c}; \texttt{src/stdio\_scan.c};
\texttt{tests/c99/numeric-text.c}; \texttt{tests/c99/run-tc-0039.ps1}.}
\def\TCTestTechnique{Boundary-value analysis using exact bit patterns,
metamorphic round trips, and adjacent-value rounding partitions.}
\def\TCPreconditions{REQ-0024 and REQ-0028 pass and the target floating model
matches the documented binary32/binary64 ABI.}
\def\TCEnvironment{C99 compilation and native execution on x86, x64, and
ARM64 with the supported TinyCC package.}
\def\TCAssumptions{The active locale is C and the default documented rounding
policy is round-to-nearest, ties-to-even.}
\def\TCInputData{Signed zeros; minimum and maximum subnormal/normal values;
halfway and adjacent values; powers of two; maximum finite; overflow,
underflow, infinity, quiet NaN, payload, malformed, case, and end-pointer
partitions.}
\def\TCInitialState{errno and output storage contain known sentinel values.}
\def\TCProcedure{Convert exact hexadecimal and special subjects through each
strto function; retain the REQ-0028 decimal regression; inspect values and end
pointers; format and scan the same bit patterns; exercise precision boundaries
and verify default-a round trips.}
\def\TCExpectedResult{Hexadecimal conversions are correctly rounded, signed
zero and classifications are retained, range errors follow the controlled
policy, ordinary success preserves errno, and finite default-a output round
trips exactly.}
\def\TCPostConditions{No persistent state other than explicitly tested errno
values remains.}
\TCTemplate
